\documentclass{article}
\usepackage[UTF8]{ctex}
\usepackage{amsmath,amssymb,amsfonts}
\usepackage{enumitem}
\usepackage{geometry}
\geometry{a4paper, margin=1in}
% === 错题本排版优化（不参与题目标识） ===
\setlength{\parindent}{0pt}
\setlength{\parskip}{0.35em plus 0.10em minus 0.08em}
\setlength{\emergencystretch}{3em}
\setlist{itemsep=0.25em, topsep=0.40em, parsep=0pt}
\allowdisplaybreaks[2]
% === 错题本排版优化结束 ===
\begin{document}
\section*{华中科技大学2026年数学分析真题}
\section*{134．华中科技大学2026年数学分析试卷（回忆版）}

\subsection{题目1}
1．（10 分）求极限 \(\lim _{n \rightarrow \infty} \sum_{k=1}^{n} \sin \frac{2 k-1}{n^{2}}\) ．

\subsection{题目2}
2．（10 分）求极限 \(\lim _{x \rightarrow 0} \frac{\tan (\tan x)-\sin (\sin x)}{\left(2026^{x}-1\right) \ln \left(1+x^{2}\right)}\) ．

\subsection{题目3}
3．（10 分）设 \(\theta \in(-1,1)\) ，试求 \(\int_{0}^{\frac{\pi}{2}} \ln \left(\frac{1+\theta \cos x}{1-\theta \cos x}\right) \frac{1}{\cos x} \mathrm{~d} x\) ．

\subsection{题目4}
4．（10 分）设 \(f(x)=\frac{1}{\sqrt{4 \pi^{2}-x^{2}}}\) 在 \([-\pi, \pi]\) 上的傅里叶级数展开式为

\[
f(x) \sim \frac{a_{0}}{2}+\sum_{n=1}^{\infty}\left(a_{n} \cos n x+b_{n} \sin n x\right) .
\]

求 \(\sum_{n=1}^{\infty}\left(a_{n}^{2}+b_{n}^{2}\right)\) ．

\subsection{题目5}
5．（10 分）设点 \(A(0,0), B\left(\pi, \pi^{\frac{2}{3}}\right), \Gamma\) 是由 \(A\) 沿 \(y=x^{\frac{2}{3}}\) 到 \(B\) 的曲线段，试求

\[
\int_{\Gamma} \frac{\cos y}{\sqrt{x^{2}+1}} \mathrm{~d} x-\ln \left(x+\sqrt{x^{2}+1}\right) \sin y \mathrm{~d} y .
\]

\subsection{题目6}
7．（15 分）设 \(x_{1}=\frac{2025}{2026}, x_{n+1}=\sin x_{n}(n=1,2, \cdots)\) ．

\begin{itemize}
  \item[（1）] 证明： \(\lim _{n \rightarrow \infty} x_{n}=0\) ．
  \item[（2）] 试证明：存在正整数 \(m\) ，使得存在 \(c \in \mathbb{R}\) ，有 \(\lim _{n \rightarrow \infty} n^{\frac{1}{m}} x_{n}=c\) ．
\end{itemize}

\subsection{题目7}
8．（15 分）\(\left\{a_{n}\right\},\left\{b_{n}\right\}\) 满足 \(\sum_{n=1}^{\infty}\left|a_{n+1}-a_{n}\right|\) 收敛， \(\lim _{n \rightarrow \infty} a_{n}=0,\left\{b_{n}\right\}\) 的部分和有界，证明：\(\sum_{n=1}^{\infty} a_{n} b_{n}\) 收敛。

\subsection{题目8}
9．（15 分）设 \(f(x)\) 一阶连续可导，且 \(f(0)=f^{\prime}(0)=0, f^{\prime \prime}(0)\) 存在，证明：

\[
\lim _{x \rightarrow 0} \frac{f(\tan x)-f(\sin x)}{x^{3} \ln (1+x)}=\frac{1}{2} f^{\prime \prime}(0) .
\]

\subsection{题目9}
10．（15 分）设 \(f(x) \in C[a, b]\) ，证明： \(\lim _{n \rightarrow \infty} \frac{1}{n} \ln \left(\int_{a}^{b} e^{-n f(x)} \mathrm{d} x\right)=-\min _{x \in[a, b]} f(x)\) ．

\subsection{题目10}
11．（15 分）证明：\(\sum_{n=1}^{\infty} \frac{\sin n x}{n} \arctan x\) 在 \([0,1]\) 上一致收敛。

\subsection{题目11}
12．（15 分）一个巨麻烦的积分不等式证明．
\end{document}
